\documentclass[11pt,a4paper]{article}

\usepackage[utf8]{inputenc}
\usepackage[T1]{fontenc}
\usepackage{lmodern}
\usepackage[margin=2.1cm]{geometry}
\usepackage{amsmath,amssymb,amsfonts}
\usepackage{booktabs}
\usepackage{graphicx}
\usepackage{float}
\usepackage{caption}
\usepackage{subcaption}
\usepackage{array}
\usepackage{xcolor}
\usepackage{listings}
\usepackage{hyperref}
\usepackage{fancyhdr}

\hypersetup{
    colorlinks=true,
    linkcolor=blue!60!black,
    citecolor=blue!60!black,
    urlcolor=blue!60!black
}

\pagestyle{fancy}
\fancyhf{}
\cfoot{\thepage}
\renewcommand{\headrulewidth}{0pt}

\captionsetup{font=small,labelfont=bf}
\setlength{\parskip}{0.25em}
\setlength{\parindent}{1.2em}
\setlength{\textfloatsep}{8pt plus 2pt minus 2pt}
\setlength{\floatsep}{6pt plus 2pt minus 2pt}
\setlength{\intextsep}{6pt plus 2pt minus 2pt}
\renewcommand{\topfraction}{0.9}
\renewcommand{\bottomfraction}{0.85}
\renewcommand{\textfraction}{0.07}
\renewcommand{\floatpagefraction}{0.75}
\setcounter{topnumber}{4}
\setcounter{bottomnumber}{4}
\setcounter{totalnumber}{6}

\lstdefinestyle{rcode}{
    language=R,
    basicstyle=\ttfamily\scriptsize,
    keywordstyle=\color{blue!70!black}\bfseries,
    commentstyle=\color{green!45!black}\itshape,
    stringstyle=\color{red!60!black},
    numbers=left,
    numberstyle=\tiny\color{gray},
    stepnumber=1,
    numbersep=6pt,
    backgroundcolor=\color{gray!5},
    frame=single,
    rulecolor=\color{gray!40},
    breaklines=true,
    breakatwhitespace=true,
    showstringspaces=false,
    tabsize=2
}

\begin{document}

\begin{titlepage}
\thispagestyle{empty}
\vspace*{-0.8cm}
\noindent
\hspace*{0.04\textwidth}%
\raisebox{-0.5\height}{\includegraphics[height=1.55cm]{logo/vgu.png}}%
\hfill
\raisebox{-0.5\height}{\includegraphics[height=1.85cm]{logo/fra-uas.png}}%
\hspace*{0.04\textwidth}

\vspace{2.2cm}
\begin{center}
    {\Large Vietnamese -- German University}\\[0.9em]
    {\Large DEPARTMENT OF COMPUTER SCIENCE}\\[2.2em]
    {\Large Frankfurt University of Applied Sciences}\\[0.9em]
    {\Large FACULTY OF COMPUTER SCIENCE AND ENGINEERING}\\[3.2em]

    {\LARGE \textbf{FINAL REPORT}}\\[1.8em]

    {\Large \textbf{WHEN AI DISAGREES WITH YOU:}}\\[0.75em]
    {\Large \textbf{HOW UNIVERSITY STUDENTS RESPOND TO}}\\[0.75em]
    {\Large \textbf{UNCERTAIN AI ANSWERS}}\\[3.2em]

    {\large
    \textbf{Supervisor:} Prof.\ Dr.\ Eva Christina Andersson\\[1.4em]
    \begin{tabular}{@{}r@{\hspace{0.8em}}l@{}}
        \textbf{Group:} & Tran Trong Nhan -- 10423160 \\[0.85em]
                        & Nguyen Minh Nhat Huy -- 10423045 \\[0.85em]
                        & Dao Hoang Minh Triet -- 10423179
    \end{tabular}
    }
\end{center}
\vfill
\end{titlepage}

\pagenumbering{roman}

% Page 2: Acknowledgment
\vspace*{1cm}
\begin{center}
    {\Large \textbf{ACKNOWLEDGMENT}}
\end{center}
\vspace{1.2em}

\noindent
We would like to extend our sincere gratitude to Prof.\ Dr.\ Eva Christina Andersson, instructor of \textit{Data Analysis in High Dimensions}, for her dedicated mentorship, constructive feedback, and steady encouragement over the course of this project. Her clear lectures and thoughtful suggestions played a vital role in guiding our methodology and improving the overall rigor of this report.

\vspace{0.8em}
\noindent
We also wish to thank every student at the Vietnamese-German University (VGU) who took part in our survey. Without their generous time and candid responses, this empirical investigation would not have been possible.

\newpage

% Page 3: Table of Contents
\tableofcontents
\newpage

% Page 4: List of Figures
\listoffigures
\newpage

% Page 5: List of Tables
\listoftables
\newpage

% Page 6: Abstract and Main Body
\pagenumbering{arabic}

\begin{abstract}
\noindent
As generative artificial intelligence (AI) assistants become widespread in university coursework, students routinely encounter uncertain or questionable outputs---such as an AI contradicting a student's own reasoning, disagreements between separate AI tools, unverifiable bibliographic citations, unclear intermediate reasoning steps, or answer reversals under simple questioning. Understanding how students react initially to such uncertainty is essential for assessing critical AI literacy in higher education. Rather than imposing predefined categories or grading student actions as right or wrong, this project seeks to uncover whether natural behavioral patterns emerge in students' immediate responses across realistic coursework conflict scenarios. Specifically, we aim to identify distinct response strategies, examine how different agglomerative hierarchical clustering linkage methods behave on nominal survey data, and evaluate whether the resulting groupings are stable and associated with students' academic seniority and frequency of AI use. To investigate these objectives, we designed and administered an anonymous scenario-based survey at the Vietnamese-German University (VGU), recording each respondent's immediate action choice across multiple AI-uncertainty scenarios alongside withheld validation variables for year of study and weekly AI usage frequency. After cleaning and harmonizing the nominal survey data, we measured pairwise student dissimilarities using the Normalized Hamming (Gower) metric to preserve the unordered categorical response space without arbitrary numerical encoding. We then constructed and compared hierarchical clusterings using Single Linkage, Complete Linkage, and Average Linkage (UPGMA), evaluated dendrogram quality and cluster separation through cophenetic correlation and silhouette analysis, tested partition robustness via row-order permutations and random subsampling with the Adjusted Rand Index, profiled the resulting behavioral clusters, and performed external statistical validation against year of study and weekly AI usage.
\end{abstract}

\noindent\textbf{Keywords:} Multivariate Statistics, Agglomerative Hierarchical Clustering, Normalized Hamming Distance, Gower Dissimilarity, UPGMA, Silhouette Analysis, Adjusted Rand Index, AI Literacy.

\newpage

\section{Introduction, Motivation, and Research Questions}

Generative AI tools such as ChatGPT, Gemini, and Claude are now a normal part of university coursework. While these tools can help students solve problems faster, they are not always reliable. Students often see AI answers that sound convincing but are wrong, include fake paper references, skip or hide math steps, or suddenly change their answer when asked a simple ``Are you sure?''. What students choose to do first when they face an unclear or conflicting AI answer decides whether AI helps them learn better or leads them to accept wrong information.

Instead of forcing students into fixed categories beforehand or grading their choices as ``right'' or ``wrong,'' our goal is to see whether natural groups of student behavior appear directly from their survey answers. Based on this goal, we focus on three main research questions:
\begin{itemize}
    \item \textbf{RQ1 (Behavioral Groups):} When an AI answer is confusing or disagrees with other information across coursework scenarios ($Q_1$--$Q_5$), do university students show clear patterns in what they choose to do first?
    \item \textbf{RQ2 (Linkage Sensitivity):} How do the resulting cluster groups change when we use Single Linkage, Complete Linkage, and Average Linkage (UPGMA) on a discrete Hamming distance space?
    \item \textbf{RQ3 (Partition Stability \& External Validity):} Do the student clusters stay stable when we shuffle the row order or take $80\%$ random subsamples, and do they differ across external variables not used in clustering---specifically Year of Study ($S_4$) and Weekly AI Usage Frequency ($Q_6$)?
\end{itemize}

\section{Theoretical Framework and Method Explanation}

\subsection{General Objectives of Hierarchical Cluster Analysis}
Following the course formulation of multivariate cluster analysis \cite{andersson2026}, cluster analysis is an unsupervised exploratory technique for partitioning a multivariate dataset of $n$ sample units measured on $p$ variables into ``natural'' homogeneous groups. We assume the $n$ observed sample units originate from an unknown number of subpopulations with no \textit{a priori} group labels. Let
\begin{equation}
    X_{ik} = \text{response for variable } k \in \{1, \dots, p\} \text{ in sample unit } i \in \{1, \dots, n\},
\end{equation}
so that each respondent $i$ is represented by a $p$-dimensional observation vector $\vec{X}_i = (X_{i1}, X_{i2}, \dots, X_{ip})^\top$.

Cluster analysis methods divide broadly into \textit{non-hierarchical methods} (such as MacQueen's $K$-means algorithm, which iteratively reallocates observations to $K$ pre-specified centroids) and \textit{hierarchical methods}. Hierarchical methods construct a nested sequence of partitions represented visually as a two-dimensional tree diagram (\textit{dendrogram}), operating in one of two directions:
\begin{enumerate}
    \item \textbf{Agglomerative Hierarchical Algorithms (Bottom-Up):} Start at step $0$ with $n$ singleton clusters $\{C_1, \dots, C_n\}$ where $C_i = \{\vec{X}_i\}$. At each subsequent step $t = 1, \dots, n-1$, the two closest existing clusters $C_r$ and $C_s$ (according to an inter-cluster linkage dissimilarity $d_{rs}$) are merged into a single cluster $C_r \cup C_s$, reducing the number of clusters by one until all $n$ sample units form a single cluster of size $n$.
    \item \textbf{Divisive Hierarchical Algorithms (Top-Down):} Start with all $n$ observations in a single cluster and recursively split one existing cluster into two daughter clusters at each step. Because divisive splitting requires evaluating $2^{n-1}-1$ possible bipartitions at the first step alone, agglomerative algorithms are used far more widely in practice.
\end{enumerate}

\subsection{Nominal Feature Space and Normalized Hamming Distance}
To cluster sample units, we first require a pairwise \textit{measure of association} (dissimilarity) $d(\vec{X}_i, \vec{X}_j)$ between observation vectors $\vec{X}_i$ and $\vec{X}_j$. As established in lecture \cite{andersson2026}, a valid metric distance function must satisfy four axiomatic properties for all $\vec{X}_i, \vec{X}_j, \vec{X}_k$:
\begin{align}
    \text{\textbf{1. Symmetry:}} \quad & d(\vec{X}_i, \vec{X}_j) = d(\vec{X}_j, \vec{X}_i) \label{eq:sym}\\
    \text{\textbf{2. Positivity:}} \quad & d(\vec{X}_i, \vec{X}_j) > 0 \quad \text{if } \vec{X}_i \neq \vec{X}_j \label{eq:pos}\\
    \text{\textbf{3. Identity:}} \quad & d(\vec{X}_i, \vec{X}_j) = 0 \quad \text{if } \vec{X}_i = \vec{X}_j \label{eq:id}\\
    \text{\textbf{4. Triangle Inequality:}} \quad & d(\vec{X}_i, \vec{X}_k) \le d(\vec{X}_i, \vec{X}_j) + d(\vec{X}_j, \vec{X}_k) \label{eq:tri}
\end{align}
Standard continuous distance metrics---such as Euclidean distance $d(\vec{X}_i, \vec{X}_j) = \sqrt{\sum_{k=1}^p (X_{ik} - X_{jk})^2}$, Minkowski distance $\left[\sum_{k=1}^p |X_{ik} - X_{jk}|^m\right]^{1/m}$, and the Canberra metric---require arithmetic differences $X_{ik} - X_{jk}$. However, in our survey, each variable $X_{ik} \in \mathcal{A}_k$ represents an \textit{unordered nominal category} (e.g., Option \texttt{A} = checking course materials, Option \texttt{C} = asking the same AI, Option \texttt{D} = asking a different AI or classmate). Encoding $\{\texttt{A}, \texttt{B}, \texttt{C}, \texttt{D}, \texttt{E}, \texttt{F}\}$ as integers $\{1, 2, 3, 4, 5, 6\}$ and computing Euclidean distances would impose an artificial linear ordering where Option \texttt{A} is five times farther from \texttt{F} than from \texttt{B}.

Instead, we define dissimilarity on the nominal space $\mathcal{X} = \mathcal{A}_1 \times \dots \times \mathcal{A}_p$ using the \textbf{Normalized Hamming Distance} (which coincides exactly with \textbf{Gower's Dissimilarity Coefficient} \cite{gower1971} for equally weighted qualitative variables with complete data):
\begin{equation}
    d_H(\vec{X}_i, \vec{X}_j) = \frac{1}{p} \sum_{k=1}^{p} \delta(X_{ik}, X_{jk}), \qquad \text{where } \delta(X_{ik}, X_{jk}) = \begin{cases} 0 & \text{if } X_{ik} = X_{jk}, \\ 1 & \text{if } X_{ik} \neq X_{jk}. \end{cases}
    \label{eq:hamming}
\end{equation}
For $p = 5$ scenarios, $d_H(\vec{X}_i, \vec{X}_j)$ measures the exact proportion of scenarios on which student $i$ and student $j$ chose different actions, taking values in the discrete six-point lattice
\begin{equation}
    d_H(\vec{X}_i, \vec{X}_j) \in \left\{0.0,\, 0.2,\, 0.4,\, 0.6,\, 0.8,\, 1.0\right\}.
\end{equation}
It is straightforward to verify that $d_H$ satisfies all four metric axioms (\ref{eq:sym})--(\ref{eq:tri}) on the space of response profiles: $\delta(a,b) = \delta(b,a)$ ensures symmetry; $\delta(a,b) \ge 0$ with equality if and only if $a=b$ ensures positivity and identity on distinct profiles $\vec{X}_i \neq \vec{X}_j$ (while two distinct students sharing an identical 5-question profile $\vec{X}_i = \vec{X}_j$ legitimately have $d_H = 0$); and since $\delta(X_{ik}, X_{kk'}) \le \delta(X_{ik}, X_{jk}) + \delta(X_{jk}, X_{kk'})$ holds coordinate-wise for the discrete metric, summing over $k=1,\dots,p$ and dividing by $p$ preserves the triangle inequality.

\subsection{Mathematical Formulation of Linkage Criteria}
Suppose cluster 1 contains $k$ observation vectors $\{\vec{X}_1, \vec{X}_2, \dots, \vec{X}_k\}$ and cluster 2 contains $l$ observation vectors $\{\vec{Y}_1, \vec{Y}_2, \dots, \vec{Y}_l\}$. Following the notation of Prof.\ Dr.\ C.\ Andersson \cite{andersson2026}, the five standard agglomerative linkage methods define inter-cluster dissimilarity $d_{12}$ as follows:

\begin{enumerate}
    \item \textbf{Single Linkage (Nearest-Neighbor):} Defines inter-cluster distance as the minimum distance between any single observation in cluster 1 and any single observation in cluster 2:
    \begin{equation}
        d_{12}^{\text{single}} = \min_{1 \le i \le k,\, 1 \le j \le l} d(\vec{X}_i, \vec{Y}_j).
    \end{equation}
    At each stage, the two clusters with the smallest $d_{12}^{\text{single}}$ are merged. While sensitive to local bridges, Single Linkage suffers from the \textit{chaining phenomenon}: two large, otherwise distant clusters are merged prematurely whenever a chain of intermediate points connects them.

    \item \textbf{Complete Linkage (Furthest-Neighbor):} Defines inter-cluster distance as the maximum distance between any single observation in cluster 1 and any single observation in cluster 2:
    \begin{equation}
        d_{12}^{\text{complete}} = \max_{1 \le i \le k,\, 1 \le j \le l} d(\vec{X}_i, \vec{Y}_j).
    \end{equation}
    Complete linkage enforces compact, spherical-diameter clusters because merging two groups requires every pair across the combined group to lie within $d_{12}^{\text{complete}}$.

    \item \textbf{Average Linkage (UPGMA --- Unweighted Pair Group Mean Averaging):} Defines inter-cluster distance as the arithmetic average of all $k \times l$ pairwise distances between observations in cluster 1 and observations in cluster 2:
    \begin{equation}
        d_{12}^{\text{average}} = \frac{1}{kl} \sum_{i=1}^{k} \sum_{j=1}^{l} d(\vec{X}_i, \vec{Y}_j).
        \label{eq:avg_link}
    \end{equation}
    Average linkage provides a compromise between the chaining tendency of Single Linkage and the extreme outlier sensitivity of Complete Linkage, serving as our primary clustering method.

    \item \textbf{Centroid Method and Ward's Minimum Variance Method (Why Not Applied Directly):}
    In the \textbf{Centroid Method}, inter-cluster distance is the distance between cluster mean vectors, $d_{12}^{\text{centroid}} = d(\bar{\vec{x}}_1, \bar{\vec{y}}_2)$. In \textbf{Ward's Method}, clustering is formulated as an ANOVA-based variance minimization problem \cite{andersson2026}, merging the pair of clusters that minimizes the increase in the combined Error Sum of Squares ($SSE$) or equivalently maximizes $R^2$:
    \begin{equation}
        SSE = \sum_{i} \sum_{j} \sum_{k=1}^{p} |X_{ijk} - \bar{x}_{i\cdot k}|^2, \qquad SST = \sum_{i} \sum_{j} \sum_{k=1}^{p} |X_{ijk} - \bar{x}_{\cdot\cdot k}|^2, \qquad R^2 = \frac{SST - SSE}{SST}.
    \end{equation}
    Because $X_{ijk}$ in our study are unordered nominal labels ($\texttt{A}$--$\texttt{F}$), sample means $\bar{x}_{i\cdot k}$ and squared deviations $|X_{ijk} - \bar{x}_{i\cdot k}|^2$ are undefined without arbitrary numerical coding or high-dimensional one-hot dummy expansion. In contrast, Single, Complete, and Average Linkage operate strictly on the pairwise dissimilarity matrix $D = [d_H(\vec{X}_i, \vec{X}_j)]_{n \times n}$, respecting the nominal scale of the survey.
\end{enumerate}

\subsection{Strengths and Weaknesses of the Methodology}
\textbf{Strengths:} (1) Agglomerative clustering with Normalized Hamming distance requires no parametric distributional assumptions or arbitrary numerical scoring of categorical choices; (2) it reveals the full multi-scale hierarchy of student behavior (e.g., showing how broad self-verification splits into textbook lookup versus independent derivation); and (3) Average Linkage smoothing transforms discrete pairwise distances into fine-grained rational cluster distances.

\textbf{Weaknesses and Theoretical Nuances:}
\begin{itemize}
    \item \textit{Equal-Distance Assumption Within Questions:} Normalized Hamming distance treats any mismatch $\delta(X_{ik}, X_{jk}) = 1$ equally,even if two options (such as checking lecture notes vs.\ solving independently) are conceptually closer than accepting AI uncritically.
    \item \textit{Discrete Distance Lattice \& Tie-Breaking Sensitivity:} With $p = 5$ questions, pairwise distances take only $6$ values. Under Single and Complete Linkage, since $\min$ and $\max$ are order statistics of $\{0, 0.2, 0.4, 0.6, 0.8, 1.0\}$, \textit{every merge height in the entire dendrogram is restricted to those 6 numbers}. Thousands of tied pairs force the algorithm to break ties according to respondent row order in the dataset. Average Linkage (\ref{eq:avg_link}) overcomes this limitation because dividing by $kl$ produces rational merge heights with far fewer ties.
    \item \textit{Varying Category Counts ($5$ vs.\ $6$ Options):} Scenarios $Q_1, Q_2, Q_5$ have $M_k = 6$ answer options ($\texttt{A}$--$\texttt{F}$), whereas $Q_3, Q_4$ have $M_k = 5$ options ($\texttt{A}$--$\texttt{E}$). While each question receives an equal weight of $1/5 = 0.20$ in Equation (\ref{eq:hamming}), the probability of two respondents matching purely by chance under uniform selection is $\frac{1}{5} = 20.0\%$ on $Q_3, Q_4$ compared to $\frac{1}{6} \approx 16.7\%$ on $Q_1, Q_2, Q_5$ (and empirically $\sum_{m} \hat{p}_{km}^2$). Thus, 6-option questions contribute slightly more to total pairwise dissimilarity on average than 5-option questions.
\end{itemize}

\section{Practical Empirical Study: Survey Design and Data Collection}

\subsection{Survey Design and the Five AI Uncertainty Scenarios}
We developed an anonymous online questionnaire via Google Forms titled \textit{``When AI Disagrees With You: How University Students Respond to Uncertain AI Answers.''} Following two screening questions verifying university enrollment ($S_1$) and voluntary informed consent ($S_2$), along with demographic items capturing study program ($S_3$) and Year of Study ($S_4$), respondents were presented with $p = 5$ realistic coursework scenarios ($Q_1$--$Q_5$) and asked to select the \textbf{single action they would most likely take first}:
\begin{itemize}
    \item \textbf{$Q_1$ (AI Disagrees With Your Answer --- 6 options):} You solve a multiple-choice question and choose B; AI chooses D with an explanation. (\texttt{A}: Check course material; \texttt{B}: Re-solve independently; \texttt{C}: Ask same AI why B is wrong; \texttt{D}: Ask a different AI; \texttt{E}: Change to D without checking; \texttt{F}: Keep B without checking).
    \item \textbf{$Q_2$ (Two AI Assistants Disagree --- 6 options):} Two AIs give plausible explanations for B vs.\ D on an unsolved problem. (\texttt{A}: Check course material; \texttt{B}: Solve independently; \texttt{C}: Send 2nd explanation to 1st AI to compare; \texttt{D}: Ask a 3rd AI; \texttt{E}: Choose preferred AI without checking; \texttt{F}: Leave unanswered and move on).
    \item \textbf{$Q_3$ (A Cited Paper Cannot Be Located --- 5 options):} AI cites a paper title that yields no result on Google Scholar. (\texttt{A}: Independently search another source to verify claim; \texttt{B}: Ask same AI for DOI/link; \texttt{C}: Ask different AI if claim is correct; \texttt{D}: Use claim but omit citation; \texttt{E}: Remove claim without investigating).
    \item \textbf{$Q_4$ (One Reasoning Step Is Unclear --- 5 options):} You cannot see why step 3 leads to step 4 in an AI solution. (\texttt{A}: Look up concept/formula in course materials; \texttt{B}: Derive step 4 from step 3 yourself; \texttt{C}: Ask same AI to explain transition; \texttt{D}: Ask a classmate; \texttt{E}: Use solution without clarifying).
    \item \textbf{$Q_5$ (AI Changes Answer Without New Evidence --- 6 options):} AI chooses B; you ask ``Are you sure?''; AI apologizes and switches to D. (\texttt{A}: Check course material; \texttt{B}: Solve independently; \texttt{C}: Ask AI to compare both and explain change; \texttt{D}: Ask different AI; \texttt{E}: Accept D without checking; \texttt{F}: Keep B without checking).
\end{itemize}
Finally, respondents answered $Q_6$: \textit{``On how many of the past seven days did you use a generative AI assistant for coursework?''} ($0$ to $7$ days). Importantly, $S_4$ and $Q_6$ were withheld from the clustering algorithm and reserved strictly for external validation.

\subsection{Data Preprocessing and Sample Demographics}
A total of $N = 263$ responses were collected. All $263$ respondents satisfied $S_1 = \text{Yes}$ and $S_2 = \text{Yes}$ with zero missing values across $Q_1$--$Q_6$. During preprocessing in \texttt{R}, respondents were assigned anonymous IDs (\texttt{R001}--\texttt{R263}). Inspection of unique text strings revealed that Option \texttt{D} in $Q_1$ appeared under two slight wording variants between the pilot and final form (\textit{``Ask a different AI the same question again.''} [$n=29$] and \textit{``Ask a different AI the original question.''} [$n=13$]); these were harmonized into a single factor level \texttt{D} ($n=42$). Across the $6 \times 6 \times 5 \times 5 \times 6 = 5{,}400$ theoretical response combinations, the $263$ students exhibited $217$ distinct 5-question profiles, with the most frequent identical profiles being \texttt{B-B-A-B-B} ($n=8$), \texttt{D-D-C-D-D} ($n=7$), \texttt{C-C-B-C-C} ($n=6$), \texttt{F-F-E-D-F} ($n=5$), and \texttt{A-A-A-A-A} ($n=4$).

By Year of Study ($S_4$), the sample comprises $63$ 1st-year ($24.0\%$), $66$ 2nd-year ($25.1\%$), $48$ 3rd-year ($18.3\%$), $68$ 4th-year ($25.9\%$), and $12$ Graduate/Master's students ($4.6\%$), with $6$ unspecified ($2.3\%$). By Major ($S_3$), Computer Science and Engineering (CSE) accounts for $52.1\%$ ($n=137$), followed by ECE ($3.8\%$), BCE ($2.7\%$), MEN ($2.7\%$), MEC ($2.3\%$), BBA ($1.9\%$), BFA ($1.9\%$), other VGU programs ($2.7\%$), and unspecified ($30.0\%$). For weekly AI usage ($Q_6$), students reported a mean of $4.10 \pm 2.12$ days/week ($\text{median} = 4.0, \text{IQR} = 3.0$).

\begin{figure}[H]
    \centering
    \includegraphics[width=0.92\textwidth]{/mnt/hdd-1tb/CSE-VGU/2026-2027/Elective-1/output/figures/fig01_marginal_responses.pdf}
    \caption{Marginal percentage distributions of student responses across scenarios $Q_1$--$Q_5$ ($n = 263$).}
    \label{fig:marginal}
\end{figure}

Figure \ref{fig:marginal} displays the marginal distributions across $Q_1$--$Q_5$. Checking course materials (\texttt{A}, $18.6\%$--$25.1\%$), solving/deriving independently (\texttt{B}, $17.9\%$--$23.6\%$), and probing the same AI (\texttt{C}, $17.9\%$--$24.3\%$) are balanced across $Q_1, Q_2, Q_4, Q_5$. In $Q_3$ (missing paper), Option \texttt{A} (\textit{independently search for another source}) reaches $40.3\%$ ($n=106$) because it combines textbook/library verification and independent literature search into a single 5-option item. Consequently, the empirical pairwise match rate $\sum_m \hat{p}_{km}^2$ is highest for $Q_3$ ($25.81\%$) and $Q_4$ ($20.69\%$) compared to the 6-option scenarios $Q_1$ ($17.70\%$), $Q_2$ ($18.12\%$), and $Q_5$ ($18.38\%$). Across all $\binom{263}{2} = 34{,}453$ respondent pairs (Figure \ref{fig:hamming}), the mean Normalized Hamming distance is $0.8016 \pm 0.2005$ ($\text{median} = 0.80$), with $104$ pairs ($0.30\%$) at $d_H = 0.0$, $559$ ($1.62\%$) at $0.2$, $2{,}184$ ($6.34\%$) at $0.4$, $6{,}220$ ($18.05\%$) at $0.6$, $12{,}422$ ($36.05\%$) at $0.8$, and $12{,}964$ ($37.63\%$) at $1.0$.

\begin{figure}[H]
    \centering
    \includegraphics[width=0.78\textwidth]{/mnt/hdd-1tb/CSE-VGU/2026-2027/Elective-1/output/figures/fig02_hamming_distribution.pdf}
    \caption{Empirical distribution of Normalized Hamming distances across all $34{,}453$ pairs of respondents.}
    \label{fig:hamming}
\end{figure}

\section{Empirical Results: Linkage Comparison and Cluster Profiling}

\subsection{Comparison of Single, Complete, and Average Linkage (RQ2)}
We constructed agglomerative dendrograms from the $263 \times 263$ Gower/Hamming matrix using \texttt{stats::hclust} under Single, Complete, and Average Linkage (Figure \ref{fig:dendro}). To quantify how faithfully each dendrogram preserves the original pairwise distances $d_H(\vec{X}_i, \vec{X}_j)$, we computed the \textbf{Cophenetic Correlation Coefficient} $r_{\text{coph}} = \text{Cor}(d_H, d_{\text{coph}})$, where $d_{\text{coph}}(i,j)$ is the dendrogram height at which respondents $i$ and $j$ first merge:
\begin{equation}
    r_{\text{coph}}^{\text{Average}} = \mathbf{0.6275} \quad > \quad r_{\text{coph}}^{\text{Complete}} = \mathbf{0.5253} \quad \gg \quad r_{\text{coph}}^{\text{Single}} = \mathbf{0.1449}.
\end{equation}
Average Linkage preserves the pairwise geometry substantially better than both alternatives. Examining the 262 merge heights reveals why Single and Complete Linkage struggle on discrete Hamming spaces:
\begin{itemize}
    \item \textbf{Single Linkage Chaining:} Because $d_{12}^{\text{single}} = \min d(\vec{X}_i, \vec{Y}_j) \in \{0, 0.2, 0.4, 0.6\}$, out of $262$ merges, $46$ occur at $h=0.0$ (identical profiles), $154$ occur at $h=0.2$ (differing on 1 question), $61$ occur at $h=0.4$, and $1$ occurs at $h=0.6$. Because almost every student shares at least 3 or 4 answers with \textit{some} neighbor, Single Linkage chains $97.3\%$--$99.6\%$ of the entire student body into one massive cluster while peeling off singletons ($258/1/1/1/1/1$ at $K=6$).
    \item \textbf{Complete Linkage Ceiling Degeneracy:} Because $d_{12}^{\text{complete}} = \max d(\vec{X}_i, \vec{Y}_j)$, merge heights are restricted to $\{0.0, 0.2, 0.4, 0.6, 0.8, 1.0\}$, and the final $15$ merges all occur at the maximum ceiling height $h = 1.0$. Consequently, cutting the Complete Linkage tree at any $K \le 16$ via \texttt{cutree} slices inside $15$ tied merges at $h = 1.0$, making the partition heavily dependent on input row order.
    \item \textbf{Average Linkage (UPGMA) Resolution:} By averaging over cluster cardinalities $k \times l$, Average Linkage produces $66$ distinct merge heights, with the top merges occurring at strictly increasing heights (0.6714, 0.6905, 0.7065, 0.7229, 0.7661, 0.7697, 0.7729, 0.7986, 0.8050, 0.8374, 0.8576, 0.8653, and 0.8732), completely eliminating tie degeneracy in the upper tree.
\end{itemize}

\begin{figure}[H]
    \centering
    \includegraphics[width=0.98\textwidth]{/mnt/hdd-1tb/CSE-VGU/2026-2027/Elective-1/output/figures/fig03_dendrograms_comparison.pdf}
    \caption{Comparison of Single, Complete, and Average Linkage dendrograms ($n = 263$) with $K=6$ cut boxes.}
    \label{fig:dendro}
\end{figure}

\subsection{Determining the Optimal Number of Clusters and Stability Analysis (RQ3)}
Table \ref{tab:linkage_eval} and Figure \ref{fig:silhouette} summarize partition quality across $K \in \{2, \dots, 8\}$ using the Average Silhouette Width $\bar{s} = \frac{1}{n}\sum_{i=1}^n \frac{b(i)-a(i)}{\max\{a(i), b(i)\}}$, negative silhouette percentage, and cluster size distribution.

\begin{table}[H]
\centering
\caption{Comparison of Silhouette Widths, Cluster Sizes, and Stability (ARI) Across Linkage Methods.}
\label{tab:linkage_eval}
\resizebox{\textwidth}{!}{%
\begin{tabular}{llccccccc}
\toprule
\textbf{Method} & $\mathbf{K}$ & \textbf{Avg Sil ($\bar{s}$)} & \textbf{Neg Sil (\%)} & \textbf{Max Clust (\%)} & \textbf{Cluster Sizes ($n_1/\dots/n_K$)} & \textbf{Perm ARI} & \textbf{Sub80 ARI} \\
\midrule
Single   & 3 & $-0.0969$ & $65.4\%$ & $99.2\%$ & \texttt{261/1/1} & $0.493 \pm 0.014$ & $0.657 \pm 0.279$ \\
Single   & 5 & $-0.1778$ & $78.7\%$ & $98.5\%$ & \texttt{259/1/1/1/1} & $0.292 \pm 0.112$ & $0.777 \pm 0.178$ \\
Single   & 6 & $-0.2113$ & $83.3\%$ & $98.1\%$ & \texttt{258/1/1/1/1/1} & $0.251 \pm 0.120$ & $0.829 \pm 0.119$ \\
\midrule
Complete & 3 & $-0.0371$ & $51.3\%$ & $96.6\%$ & \texttt{254/5/4} & $0.074 \pm 0.144$ & $0.120 \pm 0.198$ \\
Complete & 5 & $-0.0682$ & $63.5\%$ & $90.9\%$ & \texttt{239/8/7/5/4} & $0.134 \pm 0.141$ & $0.322 \pm 0.168$ \\
Complete & 6 & $-0.0810$ & $67.3\%$ & $89.0\%$ & \texttt{234/8/7/5/5/4} & $0.137 \pm 0.109$ & $0.298 \pm 0.141$ \\
\midrule
Average  & 2 & $+0.1030$ & $16.7\%$ & $93.2\%$ & \texttt{245/18} & $0.369 \pm 0.268$ & $0.294 \pm 0.260$ \\
Average  & 3 & $+0.1178$ & $17.1\%$ & $68.1\%$ & \texttt{179/66/18} & $0.504 \pm 0.159$ & $0.507 \pm 0.179$ \\
Average  & 4 & $+0.1186$ & $21.7\%$ & $68.1\%$ & \texttt{179/46/20/18} & $0.514 \pm 0.194$ & $0.477 \pm 0.208$ \\
\textbf{Average} & \textbf{5} & $\mathbf{+0.1652}$ & $18.3\%$ & $47.1\%$ & \texttt{124/55/46/20/18} & $\mathbf{0.573 \pm 0.103}$ & $\mathbf{0.555 \pm 0.131}$ \\
\textbf{Average} & \textbf{6} & $\mathbf{+0.2001}$ & $\mathbf{15.6\%}$ & $\mathbf{24.0\%}$ & \texttt{63/61/55/46/20/18} & $\mathbf{0.511 \pm 0.093}$ & $\mathbf{0.522 \pm 0.096}$ \\
Average  & 7 & $+0.1993$ & $15.2\%$ & $24.0\%$ & \texttt{63/55/53/46/20/18/8} & --- & --- \\
Average  & 8 & $+0.1985$ & $17.5\%$ & $24.0\%$ & \texttt{63/55/53/46/18/12/8/8} & --- & --- \\
\bottomrule
\end{tabular}%
}
\end{table}

Under Average Linkage, $\bar{s}$ rises monotonically from $K=2$ ($0.1030$) to a global maximum at $\mathbf{K = 6}$ ($\mathbf{\bar{s} = 0.2001}$, with the lowest proportion of negative silhouettes at $15.6\%$), before plateauing at $K=7, 8$ where small splinter clusters of size $8$ appear. Furthermore, at $K=6$ the largest cluster share drops from $47.1\%$ ($n=124$ at $K=5$) to a well-balanced $24.0\%$ ($n=63$), yielding six substantive groups of sizes $63, 61, 55, 46, 20, 18$.

\begin{figure}[H]
    \centering
    \includegraphics[width=0.92\textwidth]{/mnt/hdd-1tb/CSE-VGU/2026-2027/Elective-1/output/figures/fig04_silhouette_comparison.pdf}
    \caption{Left: Average Silhouette Width $\bar{s}$ across $K=2\dots 8$. Right: Largest cluster percentage showing severe chaining in Single Linkage vs.\ balanced partitioning in Average Linkage.}
    \label{fig:silhouette}
\end{figure}

To answer \textbf{RQ3}, we conducted two rigorous stability experiments (Figure \ref{fig:stability}):
\begin{enumerate}
    \item \textbf{Row-Order Permutation Sensitivity (100 Random Shuffles):} Because internal merges at $h \in \{0.0, 0.2, 0.4\}$ contain ties, randomly permuting respondent row order tests whether tie-breaking alters final cluster membership. Complete Linkage collapses under row shuffling ($\text{mean ARI} = 0.074$--$0.137$), whereas Average Linkage achieves substantially higher order stability, peaking at $K=5$ ($\text{ARI} = 0.573 \pm 0.103$) and $K=6$ ($\text{ARI} = 0.511 \pm 0.093$).
    \item \textbf{$80\%$ Subsampling Stability (100 Random Subsets of $n_{\text{sub}} = 210$ Without Replacement):} Comparing each subsample partition to the full-data partition via the Adjusted Rand Index confirms strong structural recovery for Average Linkage at $K=5$ ($\text{mean ARI} = 0.555, \text{median} = 0.574, \text{IQR} = 0.169$) and $K=6$ ($\text{mean ARI} = 0.522, \text{median} = 0.532, \text{IQR} = 0.108$).
\end{enumerate}

\begin{figure}[H]
    \centering
    \includegraphics[width=0.90\textwidth]{/mnt/hdd-1tb/CSE-VGU/2026-2027/Elective-1/output/figures/fig05_stability_ari.pdf}
    \caption{Left: Distribution of Adjusted Rand Index (ARI) across $100$ random $80\%$ subsamples for Average Linkage. Right: Mean row-permutation ARI across $100$ random dataset shuffles.}
    \label{fig:stability}
\end{figure}

\subsection{Behavioral Archetypes and Cluster Profiling (RQ1)}
Table \ref{tab:profiles} and Figure \ref{fig:profiles} present the within-cluster response percentages across all options of $Q_1$--$Q_5$ for the optimal $K=6$ Average Linkage partition. Notably, the hierarchy between $K=5$ and $K=6$ has a transparent substantive interpretation: at $K=5$, Clusters 3 and 4 form a combined $124$-student group of \textit{Self-Reliant \& Course-Material Verifiers} who overwhelmingly select Option \texttt{A} in $Q_3$ (\textit{independently search for another source to verify the missing paper}). At $K=6$, this super-cluster splits cleanly into \textbf{Cluster 3 (\textit{Independent Problem-Solvers}, $n=61$)} and \textbf{Cluster 4 (\textit{Course-Material Verifiers}, $n=63$)}.

\begin{table}[H]
\centering
\caption{Within-Cluster Option Percentages ($\%$) Across Scenarios $Q_1$--$Q_5$ for Average Linkage ($K = 6$). Bold entries denote the modal response exceeding $40\%$ within each cluster and scenario.}
\label{tab:profiles}
\resizebox{\textwidth}{!}{%
\begin{tabular}{llccccccc}
\toprule
\textbf{Cluster \& Behavioral Archetype} & \textbf{Scenario} & \textbf{Opt A} & \textbf{Opt B} & \textbf{Opt C} & \textbf{Opt D} & \textbf{Opt E} & \textbf{Opt F} & \textbf{Dominant Action} \\
\midrule
\textbf{C1: Cross-AI \& Peer} & $Q_1$ (Disagree) & 15.2 & 4.3 & 13.0 & \textbf{60.9} & 6.5 & 0.0 & Ask a different AI \\
\textbf{Triangulators} & $Q_2$ (2 AIs) & 6.5 & 8.7 & 13.0 & \textbf{47.8} & 15.2 & 8.7 & Ask a third AI \\
($n = 46, 17.5\%$) & $Q_3$ (No Paper) & 15.2 & 4.3 & \textbf{65.2} & 10.9 & 4.3 & --- & Ask different AI if true \\
$\bar{Q}_6 = 3.67\text{ d/wk}$ & $Q_4$ (Unclear Step) & 15.2 & 4.3 & 15.2 & \textbf{52.2} & 13.0 & --- & Ask a classmate \\
 & $Q_5$ (AI Flips) & 8.7 & 10.9 & 6.5 & \textbf{58.7} & 8.7 & 6.5 & Ask a different AI \\
\midrule
\textbf{C2: Same-AI} & $Q_1$ (Disagree) & 9.1 & 3.6 & \textbf{65.5} & 9.1 & 3.6 & 9.1 & Ask same AI why B wrong \\
\textbf{Conversationalists} & $Q_2$ (2 AIs) & 32.7 & 9.1 & \textbf{45.5} & 3.6 & 3.6 & 5.5 & Send 2nd to 1st AI \\
($n = 55, 20.9\%$) & $Q_3$ (No Paper) & 25.5 & \textbf{65.5} & 3.6 & 3.6 & 1.8 & --- & Ask same AI for DOI \\
$\bar{Q}_6 = 4.67\text{ d/wk}$ & $Q_4$ (Unclear Step) & 1.8 & 7.3 & \textbf{83.6} & 3.6 & 3.6 & --- & Ask same AI to explain \\
 & $Q_5$ (AI Flips) & 20.0 & 9.1 & \textbf{56.4} & 7.3 & 5.5 & 1.8 & Ask AI to compare both \\
\midrule
\textbf{C3: Independent} & $Q_1$ (Disagree) & 18.0 & \textbf{57.4} & 8.2 & 3.3 & 8.2 & 4.9 & Re-solve independently \\
\textbf{Problem-Solvers} & $Q_2$ (2 AIs) & 9.8 & \textbf{54.1} & 8.2 & 9.8 & 6.6 & 11.5 & Solve independently \\
($n = 61, 23.2\%$) & $Q_3$ (No Paper) & \textbf{60.7} & 18.0 & 6.6 & 3.3 & 11.5 & --- & Search another source \\
$\bar{Q}_6 = 3.89\text{ d/wk}$ & $Q_4$ (Unclear Step) & 9.8 & \textbf{70.5} & 3.3 & 3.3 & 13.1 & --- & Derive step 4 myself \\
 & $Q_5$ (AI Flips) & 6.6 & \textbf{62.3} & 13.1 & 3.3 & 14.8 & 0.0 & Solve independently \\
\midrule
\textbf{C4: Course-Material} & $Q_1$ (Disagree) & \textbf{57.1} & 7.9 & 7.9 & 11.1 & 7.9 & 7.9 & Check course material \\
\textbf{Verifiers} & $Q_2$ (2 AIs) & \textbf{54.0} & 0.0 & 9.5 & 11.1 & 14.3 & 11.1 & Check course material \\
($n = 63, 24.0\%$) & $Q_3$ (No Paper) & \textbf{68.3} & 1.6 & 6.3 & 11.1 & 12.7 & --- & Search another source \\
$\bar{Q}_6 = 4.40\text{ d/wk}$ & $Q_4$ (Unclear Step) & \textbf{47.6} & 14.3 & 7.9 & 20.6 & 9.5 & --- & Look up in course mat. \\
 & $Q_5$ (AI Flips) & \textbf{65.1} & 4.8 & 14.3 & 1.6 & 7.9 & 6.3 & Check course material \\
\midrule
\textbf{C5: Uncritical AI} & $Q_1$ (Disagree) & 0.0 & 0.0 & 10.0 & 0.0 & \textbf{85.0} & 5.0 & Change to D unverified \\
\textbf{Conformers} & $Q_2$ (2 AIs) & 10.0 & \textbf{40.0} & 10.0 & 0.0 & 35.0 & 5.0 & Solve / Pick usual AI \\
($n = 20, 7.6\%$) & $Q_3$ (No Paper) & 15.0 & 10.0 & 5.0 & \textbf{60.0} & 10.0 & --- & Use claim, omit citation \\
$\bar{Q}_6 = 2.85\text{ d/wk}$ & $Q_4$ (Unclear Step) & 25.0 & 15.0 & 10.0 & 10.0 & \textbf{40.0} & --- & Use unclarified solution \\
 & $Q_5$ (AI Flips) & 10.0 & 10.0 & 0.0 & 30.0 & \textbf{45.0} & 5.0 & Accept D unverified \\
\midrule
\textbf{C6: AI-Skeptical} & $Q_1$ (Disagree) & 0.0 & 16.7 & 0.0 & 0.0 & 5.6 & \textbf{77.8} & Keep own B unverified \\
\textbf{Anchors} & $Q_2$ (2 AIs) & 16.7 & 11.1 & 16.7 & 11.1 & 0.0 & \textbf{44.4} & Leave task unanswered \\
($n = 18, 6.8\%$) & $Q_3$ (No Paper) & 11.1 & 22.2 & 0.0 & 5.6 & \textbf{61.1} & --- & Remove claim immediately \\
$\bar{Q}_6 = 4.44\text{ d/wk}$ & $Q_4$ (Unclear Step) & 0.0 & 5.6 & 11.1 & \textbf{44.4} & 38.9 & --- & Ask classmate / Use as-is \\
 & $Q_5$ (AI Flips) & 0.0 & 5.6 & 11.1 & 5.6 & 11.1 & \textbf{66.7} & Keep B without checking \\
\bottomrule
\end{tabular}%
}
\end{table}

The six clusters represent highly coherent behavioral archetypes:
\begin{enumerate}
    \item \textbf{C4 --- Course-Material Verifiers ($n=63, 24.0\%$):} The largest group anchors verification in authoritative academic sources, checking course materials in $Q_1$ ($57.1\%$), $Q_2$ ($54.0\%$), $Q_4$ ($47.6\%$), and $Q_5$ ($65.1\%$), and searching independent scholarly literature in $Q_3$ ($68.3\%$).
    \item \textbf{C3 --- Independent Problem-Solvers ($n=61, 23.2\%$):} Rather than looking up reference answers, these students rely on their own analytical reasoning---re-solving problems independently in $Q_1$ ($57.4\%$), $Q_2$ ($54.1\%$), and $Q_5$ ($62.3\%$), deriving the missing step themselves in $Q_4$ ($70.5\%$), and independently verifying literature in $Q_3$ ($60.7\%$).
    \item \textbf{C2 --- Same-AI Conversationalists ($n=55, 20.9\%$):} These students resolve uncertainty by staying inside the chat interface and interrogating the same AI---asking why B is wrong ($65.5\%$), asking the first AI to compare explanations ($45.5\%$), requesting a DOI link from the AI ($65.5\%$), asking the AI to explain step 3$\to$4 ($83.6\%$), and asking why its answer changed ($56.4\%$).
    \item \textbf{C1 --- Cross-AI \& Peer Triangulators ($n=46, 17.5\%$):} Instead of trusting one model, they seek a second opinion from another AI ($60.9\%$ in $Q_1$, $47.8\%$ in $Q_2$, $65.2\%$ in $Q_3$, $58.7\%$ in $Q_5$) or a human classmate ($52.2\%$ in $Q_4$).
    \item \textbf{C5 --- Uncritical AI Conformers ($n=20, 7.6\%$):} A smaller but educationally critical group that defers to AI authority without verification---abandoning their own answer for D ($85.0\%$ in $Q_1$), using unverified claims while omitting the missing citation ($60.0\%$ in $Q_3$), copying unclear steps ($40.0\%$ in $Q_4$), and accepting sycophantic reversals ($45.0\%$ in $Q_5$).
    \item \textbf{C6 --- AI-Skeptical Anchors ($n=18, 6.8\%$):} The mirror opposite of C5: when AI conflicts or behaves inconsistently, they reject the AI output outright without further investigation---keeping their own answer B ($77.8\%$ in $Q_1$, $66.7\%$ in $Q_5$), dropping the unverified claim ($61.1\%$ in $Q_3$), or postponing the task ($44.4\%$ in $Q_2$).
\end{enumerate}

\begin{figure}[H]
    \centering
    \includegraphics[width=0.98\textwidth]{/mnt/hdd-1tb/CSE-VGU/2026-2027/Elective-1/output/figures/fig06_cluster_profiles.pdf}
    \caption{Within-cluster response distributions across $Q_1$--$Q_5$ for the six behavioral archetypes ($K=6$).}
    \label{fig:profiles}
\end{figure}

\subsection{External Validation with Weekly AI Usage (\texorpdfstring{$Q_6$}{Q6}) and Year of Study (\texorpdfstring{$S_4$}{S4})}
Because differences on $Q_1$--$Q_5$ are internal to how the clusters were constructed, true validation requires testing whether the clusters differ on external variables withheld from the distance calculation ($Q_6$ and $S_4$). Table \ref{tab:ext_val} and Figure \ref{fig:ext_val} present both external validation analyses.

\begin{table}[H]
\centering
\caption{External Validation of the $K=6$ Clusters Against Weekly AI Usage ($Q_6$) and Year of Study ($S_4$).}
\label{tab:ext_val}
\resizebox{\textwidth}{!}{%
\begin{tabular}{lcccccccc}
\toprule
 & \multicolumn{3}{c}{\textbf{AI Usage $Q_6$ (Days/Week)}} & \multicolumn{5}{c}{\textbf{Year of Study $S_4$ (Row $\%$ within Cluster, $n=257$)}} \\
\cmidrule(lr){2-4} \cmidrule(lr){5-9}
\textbf{Cluster ($K=6$)} & $\mathbf{n}$ & \textbf{Mean $\pm$ SD} & \textbf{Med (IQR)} & \textbf{1st Yr} & \textbf{2nd Yr} & \textbf{3rd Yr} & \textbf{4th Yr} & \textbf{Grad/Master} \\
\midrule
\textbf{C1: Cross-AI \& Peer} & 46 & $3.67 \pm 2.04$ & $3.0$ ($2.75$) & $20.0\%$ & $31.1\%$ & $15.6\%$ & $31.1\%$ & $2.2\%$ \\
\textbf{C2: Same-AI Probe} & 55 & $\mathbf{4.67 \pm 1.86}$ & $\mathbf{5.0}$ ($2.50$) & $29.6\%$ & $14.8\%$ & $27.8\%$ & $25.9\%$ & $1.9\%$ \\
\textbf{C3: Self-Solvers} & 61 & $3.89 \pm 2.29$ & $4.0$ ($4.00$) & $19.0\%$ & $\mathbf{37.9\%}$ & $10.3\%$ & $\mathbf{32.8\%}$ & $0.0\%$ \\
\textbf{C4: Course Material} & 63 & $4.40 \pm 2.01$ & $5.0$ ($3.00$) & $23.8\%$ & $23.8\%$ & $\mathbf{28.6\%}$ & $23.8\%$ & $0.0\%$ \\
\textbf{C5: AI Conformers} & 20 & $\mathbf{2.85 \pm 1.95}$ & $\mathbf{3.0}$ ($3.00$) & $\mathbf{42.1\%}$ & $5.3\%$ & $5.3\%$ & $15.8\%$ & $\mathbf{31.6\%}$ \\
\textbf{C6: Skeptic Anchors} & 18 & $4.44 \pm 2.43$ & $5.0$ ($4.75$) & $22.2\%$ & $33.3\%$ & $5.6\%$ & $16.7\%$ & $\mathbf{22.2\%}$ \\
\midrule
\multicolumn{9}{l}{\scriptsize \textbf{Statistical Tests:} $Q_6$: Kruskal-Wallis $H = 15.023, \text{df}=5, \mathbf{p = 0.0103}$; One-way ANOVA $F(5,257) = 3.148, \mathbf{p = 0.0089}$.} \\
\multicolumn{9}{l}{\scriptsize $S_4$: Pearson's $\chi^2 = 77.045$, Monte Carlo simulated $\mathbf{p = 0.0001}$ ($10{,}000$ replicates). (At $K=5$: $H$-test $p=0.0097$, $\chi^2 = 68.13, p=0.0001$).} \\
\bottomrule
\end{tabular}%
}
\end{table}

\begin{figure}[H]
    \centering
    \includegraphics[width=0.96\textwidth]{/mnt/hdd-1tb/CSE-VGU/2026-2027/Elective-1/output/figures/fig07_external_validation.pdf}
    \caption{External validation of $K=6$ clusters against Weekly AI Usage ($Q_6$, left) and Year of Study ($S_4$, right).}
    \label{fig:ext_val}
\end{figure}

Both external validation tests are statistically significant and reveal compelling behavioral insights:
\begin{itemize}
    \item \textbf{AI Usage Frequency ($Q_6$, $p = 0.0103$):} Counter to the naive assumption that heavy AI users blindly trust AI, \textbf{Uncritical AI Conformers (C5)} actually have the \textit{lowest} weekly AI usage ($\bar{Q}_6 = 2.85 \pm 1.95$ days/week, median $3.0$), whereas \textbf{Same-AI Conversationalists (C2)} have the \textit{highest} usage ($\bar{Q}_6 = 4.67 \pm 1.86$ days/week, median $5.0$), closely followed by \textbf{AI-Skeptical Anchors (C6)} ($4.44$ days) and \textbf{Course-Material Verifiers (C4)} ($4.40$ days). Frequent users have developed habits of iterative follow-up prompting or healthy skepticism toward AI hallucinations, whereas infrequent users are more prone to treat a single AI response as final authority.
    \item \textbf{Year of Study ($S_4$, $p = 0.0001$):} Active academic verification (C3 \textit{Independent Problem-Solvers} and C4 \textit{Course-Material Verifiers}) is overwhelmingly concentrated among undergraduate cohorts (Years 1--4 account for $100\%$ of known respondents in C3 and C4). Conversely, \textbf{Uncritical AI Conformers (C5)} are disproportionately composed of 1st-year freshmen ($42.1\%$, who may lack domain confidence to challenge an AI explanation) and time-constrained Graduate/Master's students ($31.6\%$).
\end{itemize}

\section{Conclusions and Limitations}

\textbf{Summary of Findings:} (1) Answering \textbf{RQ1}, VGU students do not respond uniformly to uncertain AI outputs; nearly half ($47.1\%$) engage in rigorous academic verification---either via course materials (C4, $24.0\%$) or independent mathematical derivation (C3, $23.2\%$)---while $38.4\%$ rely on secondary AI/peer consultas (C2, $20.9\%$; C1, $17.5\%$), and smaller minorities exhibit uncritical conformity (C5, $7.6\%$) or skeptical rejection (C6, $6.8\%$). (2) Answering \textbf{RQ2}, Average Linkage (UPGMA) strongly outperforms Single and Complete Linkage on low-dimensional nominal Hamming spaces by resolving discrete merge-height ties into $66$ distinct rational levels ($r_{\text{coph}} = 0.6275$). (3) Answering \textbf{RQ3}, the identified clusters exhibit substantial subsample stability ($\text{ARI} = 0.522$--$0.555$) and demonstrate strong external validity across both weekly AI usage ($p = 0.0103$) and academic seniority ($p = 0.0001$).

\textbf{Limitations:} Our findings describe self-reported first intentions among $n = 263$ VGU students rather than observed behavioral logs or academic grades. Furthermore, unweighted Hamming distance treats all nominal mismatches within a question equally and is slightly influenced by the higher baseline match probability of 5-option scenarios ($Q_3, Q_4$) relative to 6-option scenarios ($Q_1, Q_2, Q_5$). Future work could incorporate weighted semantic dissimilarities or longitudinal observational tracking.

\begin{thebibliography}{9}
\footnotesize
\bibitem{andersson2026}
Andersson, C. (2026). \textit{Data Analysis in High Dimensions: Lecture Notes on Cluster Analysis} (pp.\ 69--126). Vietnamese-German University (VGU).

\bibitem{gower1971}
Gower, J. C. (1971). A general coefficient of similarity and some of its properties. \textit{Biometrics}, 27(4), 857--871.

\bibitem{rousseeuw1987}
Rousseeuw, P. J. (1987). Silhouettes: A graphical aid to the interpretation and validation of cluster analysis. \textit{Journal of Computational and Applied Mathematics}, 20, 53--65.

\bibitem{hubert1985}
Hubert, L., \& Arabie, P. (1985). Comparing partitions. \textit{Journal of Classification}, 2(1), 193--218.
\end{thebibliography}

\newpage
\appendix
\section{Appendix: Reproducible R Code (\texttt{code/analysis.R})}
\noindent
The complete R script used for data cleaning, Gower/Hamming dissimilarity computation, hierarchical clustering, stability testing, external validation, and figure generation is included below.

\lstinputlisting[style=rcode]{/mnt/hdd-1tb/CSE-VGU/2026-2027/Elective-1/code/analysis.R}

\end{document}
